\documentclass[11pt,a4paper]{article}

\usepackage[a4paper,margin=1in]{geometry}
\usepackage{amsmath}
\usepackage[ruled,vlined,linesnumbered]{algorithm2e}
\usepackage{float}
\usepackage{graphicx}
\usepackage{booktabs}
\usepackage{hyperref}
\usepackage{subcaption}

\title{IAE Project\\[0.5em]\large Crescenzi and Takes-Kosters Algorithm}
\author{
Ayush Kanani -- 2024101049 \\
Dhyey Thummar -- 2024101024
}
\date{\today}

\begin{document}

\maketitle

\begin{abstract}
% Write a short summary of the project here.
% Mention: the graph diameter problem, the two algorithms compared,
% the datasets used, the main metrics measured, and the most important
% experimental conclusion in 4--8 sentences.
This report studies the graph diameter problem and compares two practical
diameter algorithms: Crescenzi et al.'s iFUB method and the Takes-Kosters
BoundingDiameters algorithm. Both algorithms were implemented in Python and
evaluated on one real-world Facebook graph and several synthetic graph
families, including Erd\H{o}s-R\'enyi, Barab\'asi-Albert, and Watts-Strogatz
graphs of varying sizes. We compare the algorithms using runtime, number of
breadth-first search traversals, and accuracy against an exact baseline on
smaller graphs. The goal of the study is to understand how graph structure
affects practical performance and scalability.
\end{abstract}

\section{Introduction: The Problem}

The \emph{diameter} of an undirected, unweighted graph $G=(V,E)$ is the maximum
shortest-path distance between any pair of vertices:
\[
\mathrm{diam}(G) = \max_{u,v \in V} d(u,v).
\]
Computing the diameter is a fundamental graph problem because it captures the
largest communication distance in a network and is useful in the analysis of
social networks, communication graphs, and large sparse systems.

The most direct exact approach is based on the all-pairs shortest path (APSP)
idea: compute the shortest path distance from every node to every other node,
and then take the maximum over all pairs. In an unweighted graph, this can be
done by running breadth-first search (BFS) from each vertex, giving a time cost
of $O(n(n+m))$, where $n = |V|$ and $m = |E|$. For large sparse graphs, this is
often too expensive in practice because it requires a BFS from every node.
For this reason, practical exact algorithms aim to avoid exploring every vertex
as a BFS source. In this project, we study two such methods that are designed
to work well on sparse real-world networks: Crescenzi et al.'s iFUB algorithm
and the Takes-Kosters BoundingDiameters algorithm. Our goal is to compare these
methods both conceptually and experimentally, and to understand how graph
family and graph size affect their runtime and BFS usage in practice.

\section{Algorithms Studied}

\subsection{Crescenzi: iFUB with 4-Sweep High-Degree Initialization}

The iFUB (\emph{iterative Fringe Upper Bound}) algorithm proposed by Crescenzi
et al.~\cite{crescenzi} improves practical
diameter computation by first finding a strong lower bound using BFS sweeps and
then refining upper and lower bounds layer by layer from a central candidate
vertex. In this project, the starting node for the 4-Sweep procedure is chosen
as a high-degree vertex, which often gives a stronger initial lower bound in
practice.

\paragraph{Theoretical Complexity.}
Each BFS costs $O(n+m)$. The practical efficiency of iFUB comes from requiring
only a limited number of BFS traversals on many real-world sparse graphs.
However, in the worst case, the method may still perform many BFS traversals,
leading to a worst-case running time of $O(n(n+m))$.
Its space complexity is $O(n+m)$ due to the adjacency-list graph representation
and the BFS distance/queue structures maintained during traversal.

The main intuition is to first locate a center-like starting vertex using
4-Sweep, then run BFS from that vertex and inspect its fringe layers from the
outside inward. At each layer, the algorithm computes eccentricities only where
needed and maintains lower and upper bounds on the graph diameter. Once the two
bounds coincide, the exact diameter is known.


\begin{algorithm}[H]
\caption{Crescenzi iFUB with 4-Sweep High-Degree Initialization}
\KwIn{Connected unweighted graph $G=(V,E)$}
\KwOut{Diameter of $G$}

$r_1 \leftarrow$ a highest-degree vertex of $G$\;
$a_1 \leftarrow$ farthest vertex from $r_1$ using BFS\;
$b_1 \leftarrow$ farthest vertex from $a_1$ using BFS\;
$r_2 \leftarrow$ middle vertex of a shortest path from $a_1$ to $b_1$\;
$a_2 \leftarrow$ farthest vertex from $r_2$ using BFS\;
$b_2 \leftarrow$ farthest vertex from $a_2$ using BFS\;
$u \leftarrow$ middle vertex of a shortest path from $a_2$ to $b_2$\;
$LB \leftarrow \max(\mathrm{ecc}(a_1), \mathrm{ecc}(a_2))$\;
Run BFS from $u$ and compute level sets $L_0, L_1, \dots, L_h$\;
$UB \leftarrow 2h$\;
$i \leftarrow h$\;
\While{$UB > LB$}{
    $M \leftarrow LB$\;
    \ForEach{$v \in L_i$}{
        compute $\mathrm{ecc}(v)$ by BFS if not already known\;
        $M \leftarrow \max(M, \mathrm{ecc}(v))$\;
    }
    \If{$M > 2(i-1)$}{
        \Return{$M$}\;
    }
    $LB \leftarrow M$\;
    $UB \leftarrow 2(i-1)$\;
    $i \leftarrow i-1$\;
}
\Return{$LB$}\;
\end{algorithm}

\subsection{Takes-Kosters: BoundingDiameters}

The Takes-Kosters BoundingDiameters algorithm introduced by Takes and
Kosters~\cite{takeskosters} maintains lower and upper bounds on the
eccentricities of candidate vertices and uses these bounds to prune vertices
that cannot improve the current diameter estimate. In our implementation, the
candidate selection strategy alternates between promising upper-bound and
lower-bound choices, with degree-based tie-breaking.

\paragraph{Theoretical Complexity.}
As with iFUB, each BFS costs $O(n+m)$. The algorithm can terminate much earlier
than the naive baseline when bounds tighten quickly, but in the worst case it
may still require many BFS traversals. Therefore, its worst-case running time is
also $O(n(n+m))$.
Its space complexity is also $O(n+m)$, since the implementation stores the
graph in adjacency-list form and maintains BFS state together with lower and
upper bound arrays for the vertices.

The algorithm starts with all vertices as candidates and repeatedly selects a
candidate vertex from which it runs BFS. The distances returned by that BFS are
used to tighten lower and upper bounds on the eccentricity of every candidate,
as well as global bounds on the graph diameter. Vertices that can no longer
improve the answer are pruned, which can greatly reduce the total number of BFS
traversals in practice.

\begin{algorithm}[H]
\caption{Takes-Kosters BoundingDiameters}
\KwIn{Connected unweighted graph $G=(V,E)$}
\KwOut{Diameter of $G$}

$C \leftarrow V$\;
\ForEach{$v \in V$}{
    $L[v] \leftarrow -\infty$\;
    $U[v] \leftarrow +\infty$\;
}
$D_L \leftarrow -\infty$\;
$D_U \leftarrow +\infty$\;
\While{$C \neq \emptyset$ \textbf{and} $D_L \neq D_U$}{
    choose a candidate $x \in C$ using the bound-based selection rule\;
    run BFS from $x$ and compute distances $d(x,\cdot)$\;
    $e_x \leftarrow \mathrm{ecc}(x)$\;
    $D_L \leftarrow \max(D_L, e_x)$\;
    $D_U \leftarrow \min(D_U, 2e_x)$\;
    \ForEach{$v \in C$}{
        $L[v] \leftarrow \max(L[v], \max(e_x-d(x,v), d(x,v)))$\;
        $U[v] \leftarrow \min(U[v], e_x+d(x,v))$\;
        \If{$U[v] \leq D_L$ \textbf{and} $L[v] \geq D_U/2$}{
            remove $v$ from $C$\;
        }
    }
}
\Return{$D_L$}\;
\end{algorithm}

\section{Implementation Details}

This project was implemented entirely in Python. The code is organized into
separate modules for algorithms, graph loading, graph generation, experiment
management, and plotting.

\subsection{Graph Loading and Facebook Ego-Network Merging}
The Facebook data in our repository follows the SNAP ego-network format. Each
\texttt{.edges} file corresponds to one ego network and lists connections among
the observed alter nodes of a single ego user. Our loader constructs one merged
undirected graph by adding all alter-to-alter edges, inserting the ego node for
each network, and connecting that ego node to the alters observed in its local
network. If the merged graph were disconnected, our pipeline would keep the
largest connected component for fair diameter benchmarking. In our experiments,
the resulting \texttt{facebook\_combined} graph has $4039$ nodes and $88{,}234$
edges.

\subsection{Synthetic Graph Generation}
We use three synthetic graph families to study how the algorithms react to
different structural properties. Erd\H{o}s-R\'enyi graphs represent nearly
uniform random connectivity, Barab\'asi-Albert graphs represent scale-free
networks with hubs, and Watts-Strogatz graphs represent small-world networks
with local clustering and short path lengths. For the main benchmark profile,
we generated connected graphs with $1000$, $2000$, $5000$, $10{,}000$, and
$20{,}000$ nodes for each family.

\subsection{Other Implementation Details}
Both algorithms were implemented in Python using adjacency-list graph
representations. We use a separate exact baseline only for smaller graphs,
where running BFS from every vertex is still affordable. For each benchmark
run, we record runtime, BFS traversals, and bound information. The Crescenzi
implementation follows the iFUB workflow with 4-Sweep high-degree
initialization, while the Takes-Kosters implementation follows the
BoundingDiameters bound-pruning strategy with alternating candidate selection.

\section{Experimental Setup}

\subsection{Datasets}

The experiments use one real-world graph and multiple synthetic graph families:

\begin{itemize}
    \item \textbf{Facebook Combined:} merged ego-network graph constructed from
    the provided Facebook dataset.
    \item \textbf{Erd\H{o}s-R\'enyi:} random graphs used to study behavior on
    approximately uniform random structure.
    \item \textbf{Barab\'asi-Albert:} scale-free graphs used to study behavior
    under hub-dominated structure.
    \item \textbf{Watts-Strogatz:} small-world graphs used to study behavior on
    locally clustered graphs with short path lengths.
\end{itemize}
The final report uses the \texttt{large} benchmark profile. This profile
includes synthetic graph sizes $1000$, $2000$, $5000$, $10{,}000$, and
$20{,}000$ for each family, along with the merged Facebook graph. Exact-baseline
validation was performed earlier on smaller benchmark profiles containing
$120$- and $220$-node graphs.

\subsection{Metrics Tracked}

The following metrics were recorded during experiments:

\begin{itemize}
    \item runtime in seconds,
    \item number of BFS traversals,
    \item returned diameter,
    \item exact diameter on smaller graphs,
    \item absolute error,
    \item relative error,
    \item lower bound and upper bound where relevant.
\end{itemize}
We also recorded memory usage in the benchmark CSV, although memory is treated
as a secondary metric in this report because the clearest differences between
the two algorithms appeared in runtime and BFS counts.

\subsection{Experimental Procedure}
Both algorithms were executed on exactly the same graph instances generated by
our benchmark scripts. The results were exported automatically to a CSV file,
and the plots in this report were produced directly from that CSV using a
separate plotting script. This ensured that the textual analysis and the
figures were based on the same recorded benchmark output.

\section{Results and Comparative Analysis}

\subsection{Comparison with the Exact Baseline}

The current report is based on the \texttt{large} benchmark profile, so the CSV
contains large synthetic graphs and the \texttt{facebook\_combined} graph, but
no exact-baseline rows. Therefore, this run is used mainly for scalability
analysis. In earlier smaller-profile runs, both algorithms matched the exact
diameter on all validated synthetic graphs, so the present comparison should be
read as a performance comparison between already validated implementations.

\subsection{Speed vs. Quality Tradeoff}

In the current large-profile run, the main tradeoff is computational cost, not
answer quality. The recorded lower and upper bounds coincide in the benchmark
output, so the comparison is dominated by runtime and BFS count.

On the synthetic benchmark set, Takes-Kosters was clearly faster overall.
Crescenzi required about $534.35$ seconds in total, while Takes-Kosters
required about $89.96$ seconds. The BFS counts show the same trend:
Crescenzi performed $82{,}432$ BFS traversals, whereas Takes-Kosters performed
$12{,}399$. Thus, Takes-Kosters achieved the same practical goal using far
fewer BFS explorations on large synthetic graphs.

\subsection{Effect of Graph Structure}

Graph structure had a clear effect on performance. Takes-Kosters was faster on
$4/5$ Barab\'asi-Albert graphs, $5/5$ Erd\H{o}s-R\'enyi graphs, and $5/5$
Watts-Strogatz graphs. It also used fewer BFS traversals on these families.

The Facebook graph showed the opposite behavior. On
\texttt{facebook\_combined}, both algorithms returned diameter $8$, but
Crescenzi finished in $0.029051$ seconds using only $5$ BFS traversals, while
Takes-Kosters needed $0.060778$ seconds and $10$ BFS traversals. This suggests
that Crescenzi's 4-Sweep high-degree initialization is especially effective on
this real-world graph. A likely structural reason is that ego-network-based
social graphs often contain very high-degree hub vertices that sit close to a
large fraction of the graph. Since Crescenzi begins from a highest-degree
vertex, it is naturally biased toward these hubs, which quickly yields strong
lower bounds and a center-like starting point. In contrast, more uniform
synthetic families such as Erd\H{o}s-R\'enyi graphs do not contain such extreme
hubs as consistently, so the same high-degree initialization is less helpful
and the algorithm may need many more BFS traversals before the bounds meet.

\subsection{Discussion}

The results suggest that Takes-Kosters scales more gracefully on the large
synthetic families because its bound updates prune many candidates early and
reduce the number of BFS traversals. Crescenzi, however, performs very well
when a strong initialization quickly leads to tight bounds, which appears to be
the case on the Facebook graph. Overall, Takes-Kosters is the better choice for
the large synthetic benchmarks in this project, while Crescenzi is more
effective on the tested real-world Facebook graph.

\section{Comparison Plots}

\begin{figure}[H]
\centering
\begin{subfigure}[b]{0.48\textwidth}
    \centering
    \includegraphics[width=\linewidth]{runtime_vs_nodes.png}
    \caption{Runtime vs.\ Nodes}
\end{subfigure}
\hfill
\begin{subfigure}[b]{0.48\textwidth}
    \centering
    \includegraphics[width=\linewidth]{bfs_traversals_vs_nodes.png}
    \caption{BFS Traversals vs.\ Nodes}
\end{subfigure}
\caption{Main performance comparison plots comparing runtime and BFS usage of the two algorithms across graph families.}
\label{fig:main-plots}
\end{figure}

\begin{figure}[H]
\centering
\begin{subfigure}[b]{0.48\textwidth}
    \centering
    \includegraphics[width=\linewidth]{runtime_by_family.png}
    \caption{Runtime by Family}
\end{subfigure}
\hfill
\begin{subfigure}[b]{0.48\textwidth}
    \centering
    \includegraphics[width=\linewidth]{facebook_runtime_comparison.png}
    \caption{Facebook Runtime Comparison}
\end{subfigure}
\caption{Additional comparison plots for family-wise behavior and the real-world Facebook graph.}
\label{fig:extra-plots}
\end{figure}

\section{Conclusion}

This project compared Crescenzi's iFUB algorithm and the Takes-Kosters
BoundingDiameters algorithm for graph diameter computation on one real-world
graph and several synthetic graph families. The implementations behaved
correctly in earlier validated small-graph experiments, where both algorithms
matched the exact baseline. In the large-profile experiments used for the main
comparison, the focus was therefore on scalability and practical efficiency.

The main conclusion is that Takes-Kosters is faster in practice on the large
synthetic benchmark families and also uses fewer BFS traversals. It was faster
on most Barab\'asi-Albert graphs and on all tested Erd\H{o}s-R\'enyi and
Watts-Strogatz graphs, showing better scaling as graph size increases. This
indicates that its pruning-based strategy is effective on large sparse
synthetic networks.

At the same time, the Facebook graph produced a different outcome. Both
algorithms returned diameter $8$, but Crescenzi was faster and used only half
as many BFS traversals as Takes-Kosters. This shows that graph structure plays
an important role in practical performance, and that no single algorithm is
uniformly best for every graph family.

Overall, the experiments answer the main project questions clearly: Takes-Kosters
scales better and is usually faster on the synthetic families used here, while
Crescenzi performs especially well on the tested real-world graph. Future work
could include evaluating additional real-world datasets, extending the exact
validation set, and studying how different initialization or candidate-selection
strategies change the practical behavior of both algorithms.

\begin{thebibliography}{9}

\bibitem{crescenzi}
P. Crescenzi, R. Grossi, M. Habib, L. Lanzi, and A. Marino.
\newblock \emph{On Computing the Diameter of Real-World Undirected Graphs}.
\newblock Theoretical Computer Science, 514:84--95, 2013.
\newblock \url{https://www.sciencedirect.com/science/article/pii/S0304397512008687}

\bibitem{takeskosters}
F. W. Takes and W. A. Kosters.
\newblock \emph{Determining the Diameter of Small World Networks}.
\newblock In Proceedings of the 20th ACM International Conference on
Information and Knowledge Management (CIKM 2011), pages 1191--1196, 2011.
\newblock \url{https://dl.acm.org/doi/epdf/10.1145/2063576.2063748}

\bibitem{snap}
J. Leskovec and A. Krevl.
\newblock \emph{SNAP Datasets: Stanford Large Network Dataset Collection}.
\newblock \url{https://snap.stanford.edu/data/}

\bibitem{facebookego}
J. McAuley and J. Leskovec.
\newblock \emph{Learning to Discover Social Circles in Ego Networks}.
\newblock In Advances in Neural Information Processing Systems 25 (NIPS 2012),
pages 539--547, 2012.
\newblock \url{https://snap.stanford.edu/social2012/papers/mcauley-leskovec.pdf}

\end{thebibliography}

\end{document}
